\documentclass[12pt]{article}
\usepackage{helvet}
\renewcommand{\familydefault}{\sfdefault}
\usepackage[utf8]{inputenc}
\usepackage[spanish]{babel}
\usepackage[left=2.5cm,right=2.5cm,top=3cm,bottom=3cm]{geometry}
\usepackage{parskip}
\usepackage{graphicx}
\graphicspath{ {figures/} }
\usepackage{array}
\usepackage{multirow}
\usepackage{blindtext}
\usepackage{titlesec}
\usepackage{hyperref}
\hypersetup{
    colorlinks,
    citecolor=black,
    filecolor=black,
    linkcolor=black,
    urlcolor=black
}
\ifdefined\pdfgentounicode
    \pdfgentounicode=1
\fi

\usepackage{fancyhdr}
\usepackage{lastpage}
\setlength{\headheight}{2cm}
\renewcommand{\headrulewidth}{0pt}
\renewcommand{\footrulewidth}{0pt}
\pagestyle{fancy}
\lhead{%
    \fbox{%
        \begin{minipage}[t][2cm][c]{.23\linewidth}%
            \begin{center}
                \includegraphics[width=2.5cm]{Logo.jpeg}
                CONTRATO No:
            \end{center}
        \end{minipage}%
    }% end \fbox
}% end \lhead

\chead{%
    \fbox{%
        \begin{minipage}[t][2cm][c]{.5\linewidth}%
            \begin{center}
                \textbf{Informe}\\
                \textbf{Nombre de la consultoría, proyecto,etc.}\\
                \textbf{NOMBRE DEL ESTUDIO QUE SE PRESENTA EN ESTE DOCUMENTO.}\\
            \end{center}
        \end{minipage}%
    }% end \fbox
}% end \chead

\rhead{%
    \fbox{%
        \begin{minipage}[t][2cm][c]{.23\linewidth}%
            \begin{center}
                \textbf{Revisión: 1}\\
                \vspace{5mm}
                \textbf{Página \thepage{} de \pageref{LastPage}}
            \end{center}
        \end{minipage}%
    }% end \fbox
}% end \rhead

\begin{document}

\begin{center}
\phantom{This text will be invisible} \\
\vspace{1.5cm}
Informe N°XX\\
Nombre de la asesoría, consultoría, proyecto, etc.\\
\vspace{4cm}
NOMBRE DEL ESTUDIO QUE SE PRESENTA EN ESTE DOCUMENTO.\\
\vspace{4cm}
Elaborado por:\\
Estudiante 1\\
Estudiante 2\\
Estudiante 3\\
Estudiante 4\\
\vspace{4cm}
Dependencia que desarrolla el informe\\
PROELCO\\
Barranquilla, 23 de Febrero de 2023.\\
\vspace{2cm}
\end{center}
\newpage
\phantom{This text will be invisible} \\

\begin{table}[h]
\centering
\begin{tabular}{|ccc|}
\hline
\multicolumn{3}{|c|}{ÍNDICE DE REVISIONES}                                                       \\ \hline
\multicolumn{1}{|c|}{Índice revisión} & \multicolumn{1}{c|}{Fecha modificación} & Modificaciones \\ \hline
\multicolumn{1}{|c|}{0}               & \multicolumn{1}{c|}{}                   &                \\ \hline
\multicolumn{1}{|c|}{1}               & \multicolumn{1}{c|}{}                   &                \\ \hline
\end{tabular}
\end{table}

\begin{table}[h]
\centering
\begin{tabular}{|cc|c|c|}
\hline
\multicolumn{2}{|c|}{Versión}                            & 0 & 1 \\ \hline
\multicolumn{1}{|c|}{\multirow{2}{*}{Ejecuta}}  & Nombre &   &   \\ \cline{2-4}
\multicolumn{1}{|c|}{}                          & Fecha  &   &   \\ \hline
\multicolumn{1}{|c|}{\multirow{2}{*}{Revisa}}   & Nombre &   &   \\ \cline{2-4}
\multicolumn{1}{|c|}{}                          & Fecha  &   &   \\ \hline
\multicolumn{1}{|c|}{\multirow{2}{*}{Verifica}} & Nombre &   &   \\ \cline{2-4}
\multicolumn{1}{|c|}{}                          & Fecha  &   &   \\ \hline
\multicolumn{1}{|c|}{\multirow{2}{*}{Aprueba}}  & Nombre &   &   \\ \cline{2-4}
\multicolumn{1}{|c|}{}                          & Fecha  &   &   \\ \hline
\end{tabular}
\end{table}

\newpage
\vspace*{0.1cm}

\renewcommand\contentsname{Tabla de contenido}
\tableofcontents


\newpage
\vspace*{0.1cm}

\renewcommand{\listfigurename}{Lista de figuras}
\listoffigures


\newpage
\vspace*{0.1cm}

\renewcommand{\listtablename}{Lista de tablas}
\listoftables


\newpage
\vspace*{0.1cm}

\section{Resumen ejecutivo}

\section{Objetivo}

\section{Alcance}

\section{Documentos consultados}

\section{Descripción del proyecto}

\subsection{Etapas del proyecto}

\subsection{Ubicación geográfica}

\subsection{Área de influencia}

\section{Parámetros técnicos}

\subsection{Parámetros de equipos}

\section{Estudio de flujo y cortocircuito}

\subsection{Análisis de flujo}

\subsection{Resultados de cortocircuito}

\subsection{Conclusiones y recomendaciones}


\newpage
\vspace*{0.1cm}

\begin{thebibliography}{1}
\bibliographystyle{IEEEtran}

\bibitem{ref1}
Z. Othman and A. Satria Prabuwono, “Preliminary study on Iris recognition system: Tissues of body organs in iridology,” \emph{2010 IEEE EMBS Conference on Biomedical Engineering and Sciences (IECBES)}, pp. 115–119, Nov. 2010.

\bibitem{ref2}
J. Rojas, “Herramienta Web para el apoyo al diagnóstico de enfermedades en la medicina alternativa a través de la interpretación morfológica del iris del ojo.,” Tesis, Universidad del Valle, Tulúa, 2013.

\bibitem{ref3}
S. P. Jogi and B. B. Sharma, “Methodology of Iris image analysis for clinical diagnosis,” \emph{2014 International Conference on Medical Imaging, m-Health and Emerging Communication Systems (MedCom)}, pp. 235–240, Nov. 2014.

\bibitem{ref4}
\emph{Information technology - Biometric sample quality - Part 6: Iris image data}, ISO/IEC 29794-6:2015, International Organization for Standardization, ISO, 2015. [En línea]. Disponible: https://www.iso.org/standard/54066.html

\bibitem{ref5}
M. Moncho Santonja, B. Micó-Vicent, B. Defez, J. Jordán, \& G. Peris-Fajarnes, “Hough Transform Sensitivity Factor Calculation Model Applied to the Analysis of Acne Vulgaris Skin Lesions,” \emph{Applied Sciences}, vol. 12, no. 3, pp. 1–12, Feb. 2022.

\bibitem{ref6}
I. I. Quenta, “Modelo de Diagnostico de Enfermedades Basado en la Iridolgía,” Tesis Doctoral, Universidad Mayor de San Andrés, La Paz, Bolivia, 2009.

\bibitem{ref7}
Microsoft, “Documentation for visual studio code,” \emph{Visual Studio Code}, 03-Nov-2021. [En línea]. Disponible: https://code.visualstudio.com/docs. [Acceso: 12-Nov-2022].

\bibitem{ref8}
Python Software Foundation, “What is python? Executive summary,” \emph{Python}. [En línea]. Disponible: https://www.python.org/doc/essays/blurb/. [Acceso: 12-Nov-2022].

\bibitem{ref9}
Python Software Foundation, “Tkinter - Python interface to Tcl/Tk,” \emph{Python}. [En línea]. Disponible: https://docs.python.org/3/library/tkinter.html. [Acceso: 12-Nov-2022].

\bibitem{ref10}
Dobeš, M., Martinek J., Skoupil D., Dobešová Z.,  Pospíšil J., “Human eye localization using the modified Hough transform”. \emph{Optik}, vol 117, no. 10, pp. 468-473, Elsevier 2006. ISSN 0030-4026.

\bibitem{ref11}
Dobeš M., Machala L., Tichavský P., Pospíšil J., “Human Eye Iris Recognition Using the Mutual Information”. \emph{Optik}, vol 115, no. 9, pp. 399-405, Elsevier 2004. ISSN 0030-4026.

\bibitem{ref12}
Michal Dobeš and Libor Machala, Iris Database, http://phoenix.inf.upol.cz/iris/

\end{thebibliography}

\newpage
\vspace*{0.1cm}

\section*{Anexo A}

\section*{Anexo B}

\section*{Apéndice A}

\section*{Apéndice B}

\end{document}
